% !TEX program = pdflatex
% ------------------------------------------------------------
% MATH 521 Homework Template (adapted from your MATH 421 HW10)
%
% HOW TO USE (quick):
%   1) Edit the Metadata block (HW number / due date especially).
%   2) Paste each problem statement under \problem[]{<n>}.
%   3) Write in the Scratch/Discussion box if you want.
%   4) Write your proof under "Proof." and end with \qed (or \qedhere).
%
% Notes:
% - Keeps theorem/definition boxes (tcolorbox).
% - Keeps ToC + PDF bookmarks for each problem.
% - Adds a Scratch/Discussion box per problem (optional).
% - Header bar has course + meeting info (left) and HW number (right).
% ------------------------------------------------------------

% TROUBLESHOOTING NOTE:
% If you ever see errors like `\tcb@insert@after@part` or other tcolorbox internals being
% an "Undefined control sequence", that usually means LaTeX is accidentally loading a
% DIFFERENT (often older) `tcolorbox.sty` from your project folder or local texmf tree.
% Fix: search your repo for `tcolorbox.sty` and delete/rename any local copy so TeX uses
% the TeX Live version. Also try "Clean" in LaTeX Workshop to remove stale .aux/.toc/.out.

\documentclass[12pt]{article}

% =========================
% 0) Metadata (EDIT ME)
% =========================
\newcommand{\CourseCode}{MATH 521}
\newcommand{\CourseTitle}{Analysis I}
\newcommand{\CourseSection}{LEC 002}
\newcommand{\Semester}{Spring 2026}
\newcommand{\Instructor}{Thierry Laurens}
\newcommand{\MeetingInfo}{MWF 1:20--2:10 \;|\; Van Vleck B119} % optional
\newcommand{\HomeworkNumber}{4} % <-- change each week
\newcommand{\YourName}{Alejandro De La Torre}
\newcommand{\DueDate}{Feburary 20 by 1:20 PM CDT} % <-- or set e.g. January 31, 2026

% =========================
% 1) Core math + proof tools
% =========================
\usepackage{amsmath, amssymb, amsthm}

% =========================
% 2) Lists and spacing
% =========================
\usepackage{enumitem}
\setlist[enumerate,1]{label=\textbf{(\alph*)}, leftmargin=2em, itemsep=0.25em}
\setlist[itemize]{leftmargin=1.5em, itemsep=0.25em}
\setlength{\parskip}{0.35em}
\setlength{\parindent}{0pt}

% =========================
% 3) Page geometry + header/footer
% =========================
\usepackage{geometry}
\geometry{margin=1in}

\usepackage{fancyhdr}
\pagestyle{fancy}
\fancyhf{}
\lhead{\CourseCode\ --- \CourseSection, \Semester \;(\MeetingInfo)}
\rhead{Homework \HomeworkNumber}
\cfoot{\thepage}
\renewcommand{\headrulewidth}{0.4pt}
% Fixes common fancyhdr warning about header height:
\setlength{\headheight}{15pt}
% Optional: reclaim a bit of vertical space after increasing headheight
\addtolength{\topmargin}{-2pt}

% =========================
% 4) Links (subtle)
% =========================
\usepackage[hidelinks]{hyperref}
\setcounter{tocdepth}{1}

% =========================
% 5) Quotes / figures
% =========================
\usepackage{csquotes}
\usepackage{graphicx}
\graphicspath{{images/}} % put figures in ./images/

% =========================
% 6) Simple scratch box (no tcolorbox)
% =========================
% A lightweight environment for brief planning / scratch work.
% This avoids any tcolorbox dependencies.
\newenvironment{scratch}{%
  \par\smallskip
  \noindent\textbf{Discussion \& Scratch Work}\begin{quote}\small
}{%
  \end{quote}\smallskip
}

% =========================
% 7) Lightweight theorem-like environments (optional)
% =========================
\newtheorem*{claim}{Claim}
\newtheorem*{lemma}{Lemma}
\newtheorem*{remark}{Remark}
\newtheorem{theorem}{Theorem}
\newtheorem{proposition}{Proposition}
\newtheorem{corollary}{Corollary}
\newtheorem{definition}{Definition}
\newtheorem{example}{Example}
\newtheorem{exercise}{Exercise}

% =========================
% 8) Problem header command (with TOC + PDF bookmarks)
% Usage: \problem[Optional short title]{<number>}
% =========================
\makeatletter
\newcommand{\problem}[2][]{%
  \def\prob@title{Problem #2\if\relax\detokenize{#1}\relax\else\ (\textit{#1})\fi}%
  \phantomsection%
  \pdfbookmark[1]{\prob@title}{prob#2}%
  \addcontentsline{toc}{section}{\prob@title}%
  \section*{\prob@title}%
  \vspace{-0.25em}\hrule\vspace{0.8em}%
}
\makeatother

% =========================
% 9) Title block
% =========================
\title{Homework \HomeworkNumber}
\author{\YourName \\
\CourseCode\ --- \CourseTitle \\
\CourseSection \\
Instructor: \Instructor}
\date{\DueDate}

\begin{document}
\maketitle

% ------------------ collaboration + sources ------------------
% \section*{Collaboration Statement}
% Write “I worked alone” or list collaborators / external resources.
% "I worked on this homework independently. No outside collaborators or resources were used."

\tableofcontents
\bigskip
\hrule
\bigskip

\newpage

% ============================================================
% Problems (duplicate blocks as needed)
% ============================================================

\problem[]{1}
Let $\{a_n\}_{n\ge 1}$ be a sequence.
\begin{enumerate}
\item Prove that if $\{a_n\}_{n\ge 1}$ converges to $a$, then the sequence $\{|a_n|\}_{n\ge 1}$ converges to $|a|$. (Hint: $\lvert |x| - |y| \rvert \le |x - y|$ for any $x,y\in\mathbb{R}$.)
\item Provide an example where $\{|a_n|\}_{n\ge 1}$ converges, but $\{a_n\}_{n\ge 1}$ diverges.
\end{enumerate}

\begin{scratch}
\begin{itemize}
\item \textbf{(a)} The idea is: if $a_n\to a$, then eventually $a_n$ is very close to $a$.
The absolute value function is ``distance from $0$,'' so if two numbers are close, their distances
from $0$ are also close. The key inequality is
\[
\bigl||x|-|y|\bigr|\le |x-y|,
\]
which lets us bound $\bigl||a_n|-|a|\bigr|$ by $|a_n-a|$ and then use the definition of convergence.
\item \textbf{(b)} To make $|a_n|$ converge while $a_n$ diverges, we can alternate signs:
take $a_n=(-1)^n$. Then $|a_n|=1$ for all $n$ (so it converges), but $a_n$ oscillates between $1$
and $-1$ (so it cannot converge).
\end{itemize}
\end{scratch}

\textbf{Proof.}\\
\textbf{(a)} Assume $a_n\to a$. Let $\varepsilon>0$ be given. By convergence, there exists
$N\in\mathbb{N}$ such that for all $n\ge N$ we have $|a_n-a|<\varepsilon$.
Using the inequality $\bigl||x|-|y|\bigr|\le |x-y|$ (valid for all $x,y\in\mathbb{R}$), with
$x=a_n$ and $y=a$, we obtain for all $n\ge N$:
\[
\bigl||a_n|-|a|\bigr|\le |a_n-a|<\varepsilon.
\]
Hence, by the definition of convergence, $|a_n|$ converges to $ |a|$.

\textbf{(b)} Define $a_n:=(-1)^n$. Then $|a_n|=1$ for all $n$, so $\{|a_n|\}$ converges to $1$.
However, $\{a_n\}$ does not converge: if it converged to some $L\in\mathbb{R}$, then the subsequences
$a_{2k}=1$ and $a_{2k+1}=-1$ would both have to converge to $L$, forcing $1=L=-1$, a contradiction.
Therefore $\{|a_n|\}$ converges while $\{a_n\}$ diverges. \qed

\newpage

\problem[]{2}
(Squeeze theorem). Prove that if $\{a_n\}_{n\ge 1}$, $\{b_n\}_{n\ge 1}$, and $\{c_n\}_{n\ge 1}$ are sequences such that $a_n \le b_n \le c_n$ for all $n\ge 1$ and $\lim_{n\to\infty} a_n = L = \lim_{n\to\infty} c_n$, then $\{b_n\}_{n\ge 1}$ converges and $\lim_{n\to\infty} b_n = L$.

\begin{scratch}
\begin{itemize}
\item Intuition: $b_n$ is ``trapped'' between $a_n$ and $c_n$. If both outer sequences converge to the
same limit $L$, then for any small tolerance $\varepsilon>0$, eventually both $a_n$ and $c_n$ lie inside
the interval $(L-\varepsilon,\,L+\varepsilon)$. Since $a_n\le b_n\le c_n$, the middle term must also lie
in that same interval, forcing $b_n\to L$.

\item We will use the definition of convergence:
\[
a_n\to L \iff \forall\varepsilon>0\ \exists N\in\mathbb{N}\ \text{s.t. } \forall n\ge N,\ |a_n-L|<\varepsilon,
\]
(and similarly for $c_n$). The only ``construction'' step is choosing a single index after which \emph{both}
conditions hold; we do this by taking a maximum.
\end{itemize}
\end{scratch}

\textbf{Proof.}\\
Assume that $a_n\le b_n\le c_n$ for all $n\ge 1$, and that $\lim_{n\to\infty}a_n=L=\lim_{n\to\infty}c_n$.
We prove that $b_n\to L$.

Let $\varepsilon>0$ be given. Since $a_n\to L$, there exists $N_1\in\mathbb{N}$ such that for all $n\ge N_1$,
\[
|a_n-L|<\varepsilon \quad\Longrightarrow\quad L-\varepsilon<a_n<L+\varepsilon.
\]
Similarly, since $c_n\to L$, there exists $N_2\in\mathbb{N}$ such that for all $n\ge N_2$,
\[
|c_n-L|<\varepsilon \quad\Longrightarrow\quad L-\varepsilon<c_n<L+\varepsilon.
\]
Set $N:=\max\{N_1,N_2\}$. Then for all $n\ge N$, both inequalities hold, and we have
\[
L-\varepsilon<a_n\le b_n\le c_n<L+\varepsilon.
\]
In particular, for all $n\ge N$,
\[
L-\varepsilon<b_n<L+\varepsilon,
\]
which is equivalent to $|b_n-L|<\varepsilon$. Since $\varepsilon>0$ was arbitrary, this proves that $b_n\to L$.
\qed

\newpage

\problem[]{3}
Suppose that the sequence $\{a_n\}_{n\ge 1}$ converges.
\begin{enumerate}
\item Prove that if $a_n \ge m$ for all but finitely many $n$, then $\lim_{n\to\infty} a_n \ge m$.
\item Prove that if $a_n \le M$ for all but finitely many $n$, then $\lim_{n\to\infty} a_n \le M$.
\end{enumerate}

\begin{scratch}
\begin{itemize}
\item The key idea is: convergence prevents the sequence from staying \emph{strictly below} a number that it is eventually always above. If $a_n\to \ell$ and (eventually) $a_n\ge m$, then the limit \emph{cannot} satisfy $\ell<m$, because choosing $\varepsilon=\frac{m-\ell}{2}$ would force $a_n<\ell+\varepsilon<m$ for all large $n$, contradicting $a_n\ge m$ for all large $n$.
\item Part (b) is the symmetric argument: if (eventually) $a_n\le M$, then the limit cannot be $>M$.
\end{itemize}
\end{scratch}

\textbf{Proof.}\\
Let $\ell:=\lim_{n\to\infty} a_n$.

\textbf{(a)} Assume that $a_n\ge m$ for all but finitely many $n$. Then there exists $N_0\in\mathbb{N}$ such that
\[
 a_n\ge m\quad\text{for all }n\ge N_0.
\]
We claim that $\ell\ge m$. Suppose for contradiction that $\ell<m$. Set
\[
\varepsilon:=\frac{m-\ell}{2}>0.
\]
Since $a_n\to \ell$, there exists $N_1\in\mathbb{N}$ such that for all $n\ge N_1$,
\[
|a_n-\ell|<\varepsilon\quad\Longrightarrow\quad a_n<\ell+\varepsilon=\frac{\ell+m}{2}<m.
\]
Let $N:=\max\{N_0,N_1\}$. Then for all $n\ge N$ we have simultaneously $a_n\ge m$ and $a_n<m$, a contradiction. Hence $\ell\ge m$.

\textbf{(b)} Assume that $a_n\le M$ for all but finitely many $n$. Then there exists $N_0\in\mathbb{N}$ such that $a_n\le M$ for all $n\ge N_0$. We claim that $\ell\le M$. Suppose for contradiction that $\ell>M$. Set
\
\varepsilon:=\frac{\ell-M}{2}>0.
\]
Since $a_n\to \ell$, there exists $N_1\in\mathbb{N}$ such that for all $n\ge N_1$,
\[
|a_n-\ell|<\varepsilon\quad\Longrightarrow\quad a_n>\ell-\varepsilon=\frac{\ell+M}{2}>M.
\]
Let $N:=\max\{N_0,N_1\}$. Then for all $n\ge N$ we have simultaneously $a_n\le M$ and $a_n>M$, a contradiction. Hence $\ell\le M$.\qed

\newpage

\problem[]{4}
Suppose that the sequence $\{a_n\}_{n\ge 1}$ converges. Prove or disprove the following statements:
\begin{enumerate}
\item If $\lim_{n\to\infty} a_n \ge m$, then there exists $N\in\mathbb{N}$ so that $a_n \ge m$ for all $n\ge N$.
\item If $\lim_{n\to\infty} a_n > m$, then there exists $N\in\mathbb{N}$ so that $a_n > m$ for all $n\ge N$.
\end{enumerate}

\begin{scratch}
\begin{itemize}
\item \textbf{Key idea.} Part (a) smells suspicious because the hypothesis allows equality:
$\lim a_n \ge m$ includes the borderline case $\lim a_n = m$. Convergence to $m$ does \emph{not} force
the tail of the sequence to lie on one side of $m$; it only forces the tail to lie \emph{near} $m$.
So we should try to build a convergent sequence with limit $m$ that stays strictly below $m$ forever.

\item Part (b) replaces $\ge$ by $>$, which creates a positive gap $L-m$. If $a_n\to L$ and $L>m$,
then choosing $\varepsilon$ smaller than this gap forces $a_n$ to eventually lie in $(L-\varepsilon,L+\varepsilon)$,
which sits entirely above $m$. So (b) should be true, and the proof will be a direct $\varepsilon$--$N$ argument.
\end{itemize}
\end{scratch}

\textbf{Proof.}\\
Let $\ell:=\lim_{n\to\infty} a_n$.

\textbf{(a) Disprove.} The statement is false. Fix $m\in\mathbb{R}$ and define
\[
a_n:=m-\frac{1}{n}\qquad (n\ge 1).
\]
Then $a_n\to m$ as $n\to\infty$, so $\lim_{n\to\infty}a_n=m\ge m$ holds. However, for every $n\ge 1$ we have
$a_n=m-\frac{1}{n}<m$, so there is \emph{no} $N\in\mathbb{N}$ such that $a_n\ge m$ for all $n\ge N$.
Hence (a) is false.

\textbf{(b) Prove.} Assume $\ell>m$. Set
\[
\varepsilon:=\frac{\ell-m}{2}>0.
\]
Since $a_n\to \ell$, there exists $N\in\mathbb{N}$ such that for all $n\ge N$,
\[
|a_n-\ell|<\varepsilon \quad\Longrightarrow\quad a_n>\ell-\varepsilon
= \ell-\frac{\ell-m}{2}=\frac{\ell+m}{2}>m.
\]
Therefore $a_n>m$ for all $n\ge N$, as desired. \qed

\newpage

\problem[]{5}
Prove that if $\{a_n\}_{n\ge 1}$ is a sequence of positive numbers that converges to $a > 0$, then the sequence $\{\sqrt{a_n}\}_{n\ge 1}$ converges to $\sqrt{a}$.

\begin{scratch}
\begin{itemize}
\item The main trick is to rewrite the difference of square-roots by “rationalizing”:
\[
|\sqrt{x}-\sqrt{y}|=\frac{|x-y|}{\sqrt{x}+\sqrt{y}}\qquad (x,y\ge 0).
\]
Here we will take $x=a_n$ and $y=a$.

\item Because $a>0$, the denominator $\sqrt{a_n}+\sqrt{a}$ is always at least $\sqrt{a}$ (since $a_n>0\Rightarrow \sqrt{a_n}\ge 0$). So we can control $|\sqrt{a_n}-\sqrt{a}|$ directly by controlling $|a_n-a|$.

\item Concretely, given $\varepsilon>0$, it suffices to make $|a_n-a|<\varepsilon\sqrt{a}$.
\end{itemize}
\end{scratch}

\textbf{Proof.}\\
Let $\varepsilon>0$ be given. Since $a_n\to a$, there exists $N\in\mathbb{N}$ such that for all $n\ge N$,
\[
|a_n-a|<\varepsilon\sqrt{a}.
\]
For $n\ge N$, we have $a_n>0$ and $a>0$, so $\sqrt{a_n}$ and $\sqrt{a}$ are defined and nonnegative. Using
\[
(\sqrt{a_n}-\sqrt{a})(\sqrt{a_n}+\sqrt{a})=a_n-a,
\]
we obtain
\[
|\sqrt{a_n}-\sqrt{a}|=\frac{|a_n-a|}{\sqrt{a_n}+\sqrt{a}}\le \frac{|a_n-a|}{\sqrt{a}}.
\]
Therefore, for all $n\ge N$,
\[
|\sqrt{a_n}-\sqrt{a}|\le \frac{|a_n-a|}{\sqrt{a}}<\frac{\varepsilon\sqrt{a}}{\sqrt{a}}=\varepsilon.
\]
Since $\varepsilon>0$ was arbitrary, it follows that $\sqrt{a_n}\to \sqrt{a}$. \qed

\newpage

\problem[]{6}
Consider the following sequence: $a_1 = \sqrt{2}$ and $a_{n+1} = \sqrt{2 + a_n}$ for all $n\ge 1$.
\begin{enumerate}
\item Prove that the sequence $\{a_n\}_{n\ge 1}$ is bounded.
\item Prove that the sequence is increasing.
\item Conclude that the sequence converges, and find its limit.
\end{enumerate}

\begin{scratch}
\begin{itemize}
\item \textbf{What we need.} To show the sequence is \emph{bounded}, we need to find real numbers $m,M$ such that
$m\le a_n\le M$ for all $n\ge 1$ (equivalently: bounded below and bounded above). To show it is \emph{increasing}, we need $a_n\le a_{n+1}$ for all $n\ge 1$.

\item \textbf{Boundedness idea.} Because $a_{n+1}=\sqrt{2+a_n}$, a natural guess is that the sequence stays below $2$:
if $a_n\le 2$, then $a_{n+1}=\sqrt{2+a_n}\le \sqrt{4}=2$. Also $a_n\ge 0$ implies $a_{n+1}\ge 0$. Since $a_1=\sqrt2\in(0,2)$, induction should give $0<a_n\le 2$ for all $n$.

\item \textbf{Monotonicity idea.} Since all terms are positive, we can compare $a_{n+1}$ and $a_n$ by squaring:
\[
 a_{n+1}\ge a_n \iff \sqrt{2+a_n}\ge a_n \iff 2+a_n\ge a_n^2.
\]
The last inequality is $(a_n-2)(a_n+1)\le 0$, which is true whenever $-1\le a_n\le 2$. Our boundedness work gives $0<a_n\le 2$, so this will force $a_{n+1}\ge a_n$.

\item \textbf{Convergence + limit.} Once we know $\{a_n\}$ is increasing and bounded, we may invoke the Monotone Convergence Theorem (proved in lecture): the sequence converges. If $a_n\to L$, then passing to the limit in the recursion gives
$L=\sqrt{2+L}$, so $L^2-L-2=0$ and $L\in\{2,-1\}$. Since $L\ge 0$, we conclude $L=2$.
\end{itemize}
\end{scratch}

\textbf{Proof.}\\
\textbf{(a) Boundedness.} We first show that $0<a_n\le 2$ for all $n\ge 1$.

\emph{Base case.} We have $a_1=\sqrt2$, hence $0<a_1\le 2$.

\emph{Inductive step.} Assume $0<a_n\le 2$ for some $n\ge 1$. Then $2+a_n\le 4$, so
\[
 a_{n+1}=\sqrt{2+a_n}\le \sqrt4=2.
\]
Also $2+a_n>0$, so $a_{n+1}=\sqrt{2+a_n}>0$. Thus $0<a_{n+1}\le 2$.

By induction, $0<a_n\le 2$ for all $n\ge 1$. In particular, the sequence is bounded (below by $0$ and above by $2$).

\textbf{(b) Increasing.} Fix $n\ge 1$. From part (a), $a_n>0$, so we may square inequalities without changing direction. We compute:
\[
 a_{n+1}\ge a_n
 \iff \sqrt{2+a_n}\ge a_n
 \iff 2+a_n\ge a_n^2
 \iff a_n^2-a_n-2\le 0
 \iff (a_n-2)(a_n+1)\le 0.
\]
But part (a) gives $0<a_n\le 2$, hence $a_n-2\le 0$ and $a_n+1>0$, so indeed $(a_n-2)(a_n+1)\le 0$. Therefore $a_{n+1}\ge a_n$ for all $n\ge 1$, and the sequence is increasing.

\textbf{(c) Convergence and limit.} By parts (a) and (b), $\{a_n\}$ is bounded and increasing, so by the Monotone Convergence Theorem it converges. Let
\[L:=\lim_{n\to\infty} a_n.\]
Because $a_{n+1}=\sqrt{2+a_n}$ and the square-root function is continuous on $(0,\infty)$, we may take limits on both sides to obtain
\[
L=\lim_{n\to\infty} a_{n+1}=\lim_{n\to\infty}\sqrt{2+a_n}=\sqrt{2+\lim_{n\to\infty}a_n}=\sqrt{2+L}.
\]
Squaring gives $L^2=2+L$, i.e.
\[
L^2-L-2=0 \quad\Longrightarrow\quad (L-2)(L+1)=0.
\]
Thus $L\in\{2,-1\}$. Since each $a_n>0$, we have $L\ge 0$, so $L\ne -1$ and therefore $L=2$.
\qed

% ============================================================
% Appendix for Lectures 10-12 summarizing the key results and concepts we proved in class.
% ============================================================
\newpage
\appendix
\section*{Appendix: Toolbox}
\addcontentsline{toc}{section}{Appendix: Toolbox}

\subsection*{Completeness / Real numbers}

\begin{proposition}[Archimedean property]
For every $x\in\mathbb{R}$, there exists $n\in\mathbb{N}$ such that $n>x$.
\end{proposition}

\begin{corollary}
For every $\varepsilon>0$, there exists $n\in\mathbb{N}$ such that $\frac{1}{n}<\varepsilon$.
\end{corollary}

\begin{lemma}[Integer ``sandwich'']
For every $a\in\mathbb{R}$, there exists $N\in\mathbb{Z}$ such that
\[
N\le a < N+1.
\]
\end{lemma}

\begin{proposition}[$\mathbb{Q}$ is dense in $\mathbb{R}$]
If $x,y\in\mathbb{R}$ satisfy $x<y$, then there exists $q\in\mathbb{Q}$ such that
\[
x<q<y.
\]
\end{proposition}

\begin{corollary}[$\mathbb{R}\setminus\mathbb{Q}$ is dense in $\mathbb{R}$]
If $x,y\in\mathbb{R}$ satisfy $x<y$, then there exists $r\in\mathbb{R}\setminus\mathbb{Q}$ such that
\[
x<r<y.
\]
\end{corollary}

\subsection*{Countability}

\begin{proposition}[Uncountability via Cantor diagonal]
Let
\[
A:=\{f:\mathbb{N}\to\{0,1\}\}.
\]
Then $A$ is uncountable.
\end{proposition}

\begin{remark}[Cantor's diagonal argument (what is happening conceptually)]
If we (incorrectly) assume $A$ is countable, we can list its elements as
\[
A=\{f_1,f_2,f_3,\dots\},\qquad f_k:\mathbb{N}\to\{0,1\}.
\]
Think of this as an infinite table whose $(k,n)$-entry is $f_k(n)\in\{0,1\}$.  We then \emph{define a new function}
\[
g:\mathbb{N}\to\{0,1\},\qquad g(n):=1-f_n(n).
\]
This flips the diagonal entries.  For each $n\in\mathbb{N}$ we have
\[
g(n)\ne f_n(n),
\]
so $g\ne f_n$ for every $n$ (they differ at input $n$).  Thus $g\in A$ but $g$ is not in our purported enumeration, contradicting that the list contained all of $A$.
\end{remark}

\begin{corollary}
The interval $(0,1)=\{x\in\mathbb{R}:0<x<1\}$ is uncountable.
\end{corollary}

\subsection*{Sequences and limits}

\begin{definition}[Sequence]
A \emph{sequence of real numbers} is a function $f:\mathbb{N}\to\mathbb{R}$.  We write $a_n:=f(n)$ and denote the sequence by
\[
(a_n)_{n\ge 1}\quad\text{or}\quad \{a_n\}_{n\ge 1}.
\]
\end{definition}

\begin{definition}[Absolute value]
For $x\in\mathbb{R}$,
\[
|x|=
\begin{cases}
 x & \text{if }x\ge 0,\\
 -x & \text{if }x<0.
\end{cases}
\]
\end{definition}

\begin{exercise}[Basic properties of $|\cdot|$]
For all $x,y\in\mathbb{R}$:
\begin{enumerate}
\item $|x|\ge 0$.
\item $|x|=0\iff x=0$.
\item $|x+y|\le |x|+|y|$ \;(triangle inequality).
\item $\bigl||x|-|y|\bigr|\le |x-y|$.
\item $|xy|=|x||y|$.
\end{enumerate}
\end{exercise}

\begin{definition}[Limit / convergence]
We say that the sequence $\{a_n\}_{n\ge 1}$ \emph{converges} to $a\in\mathbb{R}$ if
\[
\forall\varepsilon>0\ \exists N\in\mathbb{N}\ \text{s.t.}\ \forall n\ge N,\ |a_n-a|<\varepsilon.
\]
In this case we write $\lim_{n\to\infty} a_n=a$ or $a_n\to a$ as $n\to\infty$.
\end{definition}

\begin{definition}[Divergence]
A sequence \emph{diverges} if it does not converge.
\end{definition}

\begin{lemma}[Uniqueness of limits]
If $\{a_n\}_{n\ge1}$ converges, then its limit is unique.
\end{lemma}

\begin{lemma}[Convergent $\Rightarrow$ bounded]
If $\{a_n\}_{n\ge1}$ converges, then it is bounded.
\end{lemma}

\begin{definition}[Bounded sequence]
We say $\{a_n\}_{n\ge1}$ is bounded above / bounded below / bounded if the set $\{a_n:n\ge1\}$ is bounded above / bounded below / bounded in $\mathbb{R}$.
\end{definition}

\begin{theorem}[Scalar multiple]
If $a_n\to a$ in $\mathbb{R}$ and $k\in\mathbb{R}$, then $(ka_n)\to ka$.
\end{theorem}

\begin{theorem}[Limit laws]
If $a_n\to a$ and $b_n\to b$, then:
\begin{enumerate}
\item $a_n+b_n\to a+b$.
\item $a_n b_n\to ab$.
\item If $b\ne 0$ and $b_n\ne 0$ for all $n$, then $\displaystyle \frac{a_n}{b_n}\to \frac{a}{b}$.
\end{enumerate}
\end{theorem}

\begin{definition}[Monotone sequences]
A sequence $\{a_n\}_{n\ge1}$ is
\begin{enumerate}
\item \emph{increasing} if $a_n\le a_{n+1}$ for all $n\ge1$,
\item \emph{strictly increasing} if $a_n<a_{n+1}$ for all $n\ge1$,
\item \emph{decreasing} if $a_n\ge a_{n+1}$ for all $n\ge1$,
\item \emph{strictly decreasing} if $a_n>a_{n+1}$ for all $n\ge1$,
\item \emph{monotone} if it is increasing or decreasing.
\end{enumerate}
\end{definition}

\begin{theorem}[Monotone convergence theorem]
If $\{a_n\}_{n\ge1}$ is monotone and bounded, then it converges.
\end{theorem}

\subsection*{Cauchy sequences and subsequences}

\begin{definition}[Cauchy sequence]
A sequence $\{a_n\}_{n\ge1}$ is \emph{Cauchy} if
\[
\forall\varepsilon>0\ \exists N\in\mathbb{N}\ \text{s.t.}\ \forall n,m\ge N,\ |a_n-a_m|<\varepsilon.
\]
\end{definition}

\begin{theorem}[Cauchy criterion]
A sequence of real numbers is Cauchy if and only if it converges.
\end{theorem}

\begin{lemma}[Cauchy $\Rightarrow$ bounded]
If $\{a_n\}_{n\ge1}$ is Cauchy, then it is bounded.
\end{lemma}

\begin{definition}[Subsequence]
Let $\{a_n\}_{n\ge1}$ be a sequence.  We say $\{a_{h_n}\}_{n\ge1}$ is a \emph{subsequence} of $\{a_n\}_{n\ge1}$ if $\{h_n\}_{n\ge1}$ is a strictly increasing sequence of natural numbers (i.e. $h_n\in\mathbb{N}$ and $h_n<h_{n+1}$ for all $n$).
\end{definition}

\begin{lemma}
If $\{h_n\}_{n\ge1}$ is strictly increasing in $\mathbb{N}$, then $h_n\ge n$ for all $n\ge1$.
\end{lemma}

\begin{example}[Subsequences]
\begin{enumerate}
\item If $h_n=n$, then $\{a_{h_n}\}$ is the original sequence.
\item If $h_n=2n$, then $\{a_{2n}\}$ is the even-index subsequence.
\item If $a_n=(-1)^n$, then the subsequence $a_{2n}=1$ converges (even though $a_n$ diverges).
\end{enumerate}
\end{example}

\begin{theorem}[Subsequence of a convergent sequence]
If $a_n\to a$ and $\{a_{h_n}\}$ is any subsequence, then $a_{h_n}\to a$.
\end{theorem}

\begin{theorem}[Convergence detected by subsequences]
A sequence $\{a_n\}$ converges to $a$ if and only if every subsequence of $\{a_n\}$ converges to $a$.
\end{theorem}

\end{document}
